% =====================================================================
% Paper A
% A pseudo-Brewster compensation collapses Fresnel dosimetry to geometry:
%   a closed-form local absorption law for the human body at mmWave
% Author: Robin Wydaeghe
% Target: IEEE Transactions on Antennas and Propagation
% =====================================================================

\documentclass[journal,twocolumn,10pt]{IEEEtran}

\usepackage[utf8]{inputenc}
\usepackage[T1]{fontenc}
\usepackage{lmodern}
\usepackage{microtype}

\usepackage{amsmath,amssymb,amsthm,mathtools}
\usepackage{bm}

\usepackage{graphicx}
\graphicspath{{../../theory/figures/}{../../validation/}{figures/}}
\usepackage{booktabs}
\usepackage{array}
\usepackage{caption}
\usepackage{subcaption}

\usepackage{cite}
\usepackage{xcolor}
\usepackage[colorlinks=true,linkcolor=blue,citecolor=blue,urlcolor=blue]{hyperref}
\usepackage[capitalize]{cleveref}

\theoremstyle{plain}
\newtheorem{theorem}{Theorem}
\newtheorem{proposition}[theorem]{Proposition}
\newtheorem{corollary}[theorem]{Corollary}
\theoremstyle{remark}
\newtheorem{remark}[theorem]{Remark}

\crefname{proposition}{Proposition}{Propositions}
\Crefname{proposition}{Proposition}{Propositions}
\crefname{corollary}{Corollary}{Corollaries}
\Crefname{corollary}{Corollary}{Corollaries}
\crefname{remark}{Remark}{Remarks}
\Crefname{remark}{Remark}{Remarks}

% --- math macros ---
\newcommand{\khat}{\hat{\bm{k}}}
\newcommand{\nhat}{\hat{\bm{n}}}
\newcommand{\rr}{\mathbf{r}}
\newcommand{\EE}{\mathbf{E}}
\newcommand{\Sinc}{S_{\mathrm{inc}}}
\newcommand{\Sab}{S_{\mathrm{ab}}}
\newcommand{\Teff}{T_{\mathrm{eff}}}
\newcommand{\Tavg}{T_{\mathrm{avg}}}
\newcommand{\Tbar}{\bar{T}}
\newcommand{\ntilde}{\tilde{n}}
\newcommand{\diff}{\mathrm{d}}
\newcommand{\pospart}[1]{\left(#1\right)_{+}}
\newcommand{\Vis}{V}
\DeclareMathOperator{\ReLU}{ReLU}
\DeclareMathOperator{\RE}{Re}

\begin{document}

\title{Pseudo-Brewster compensation collapses Fresnel
dosimetry to geometry: a closed-form local absorption
law for the human body at millimetre waves}

\author{Robin~Wydaeghe%
  \thanks{R. Wydaeghe is with Ghent University and imec, Ghent,
  Belgium (e-mail: robin.wydaeghe@ugent.be).}}

\markboth{IEEE Transactions on Antennas and Propagation,
Vol.~XX, No.~X, Month~2026}%
{Wydaeghe: Pseudo-Brewster dosimetry}

\maketitle

\begin{abstract}
Millimetre-wave wireless systems above 6\,GHz illuminate the human body
with fields whose absorption regulators specify on the body surface,
yet full-wave simulation requires meshes of order $10^{12}$ cells per
configuration. We derive a closed-form local absorption law,
$\Sab(\rr) = \Sinc\,T_0\,\Vis(\rr,\khat)\,\pospart{\nhat\cdot(-\khat)}$,
valid for any biological tissue at any frequency in the 10 to 60\,GHz
band where the body is opaque. Here $(\cdot)_+$ is the positive-part
operator and $\Vis(\rr,\khat) \in \{0,1\}$ is the binary self-shadowing
visibility, equal to one where the surface point $\rr$ is reached by an
unobstructed ray from direction $\khat$ and zero where it is shadowed
by another body part. The collapse rests on two ingredients. A
polarisation-aware treatment of the surface flux gives an exact
identity in which incidence direction and surface normal enter only
through the positive-part cosine $\pospart{\nhat\cdot(-\khat)}$. The
Fresnel transmission averaged over polarisation is nearly
angle-independent for biological tissue, a TM-rising and TE-falling
compensation that places skin within $5.6\,\%$ of its
normal-incidence value $T_0$ over $[0^\circ, 75^\circ]$ at $28\,$GHz.
The remaining geometric factor depends on body shape alone. We
validate the law against exact Mie theory for lossy spheres (errors
below $5\,\%$ above $60\,$GHz, $10\,\%$ at body-part sizes in the
mmWave band), against full Fresnel integration on the Thelonious
anatomical phantom ($0.35\,\%$ on total power), and against full
Sim4Life FDTD on Thelonious at $7\,$GHz on the IEC/IEEE\,63195 peak
$4\,\mathrm{cm}^2$ averaged metric ($1.027$ direction-averaged
ratio). For multiple incident waves, each path carries its own visibility
$\Vis(\rr,\khat_n)$ and the formula has the structure of a
single-hidden-layer rectified network whose weights are the path
directions and powers from a ray tracer. The visibility factor is the
ambient-occlusion primitive of computer graphics, evaluated at frame
rate on consumer GPUs. The chain from ray tracer to dosimetry is
therefore differentiable end-to-end with the same arithmetic primitives
as a hardware-accelerated rendering pipeline.
\end{abstract}

\begin{IEEEkeywords}
Absorbed power density, ambient occlusion, Fresnel transmission,
biological tissue, millimetre wave, pseudo-Brewster angle, ICNIRP
compliance, differentiable ray tracing.
\end{IEEEkeywords}

\IEEEpeerreviewmaketitle


\section{Introduction}\label{sec:introduction}
\IEEEPARstart{M}{illimetre-wave} wireless systems above $6\,$GHz are
central to current and future networks, and regulatory compliance
above this frequency is set locally rather than as a whole-body
average. ICNIRP 2020 and IEC/IEEE\,63195 require evaluation of the
absorbed power density $\Sab$ on the body surface
\cite{ICNIRP2020}, defined as the electromagnetic power deposited
per unit area in the superficial tissue layer. At these frequencies
the skin depth is $\delta \sim 0.3$ to $0.5\,$mm, so the penetrating
power is absorbed within less than a millimetre of the surface.
Direct simulation of $\Sab$ by the finite-difference time-domain
method (FDTD) requires a mesh fine enough to resolve this thin
absorption layer in a body that spans approximately one metre. The
resulting cell count is of order $10^{12}$ per plane-wave
configuration. A typical regulatory campaign spans tens of
frequencies, several body postures, and many incidence directions,
and runs for weeks on GPU clusters.

The literature has converged on a partial closed form. Kodera, Diao,
Funahashi, and Li \cite{Kodera2024,Diao2024,Funahashi2018,Li2019}
report empirically that whole-body absorbed power above
$10\,$GHz is well approximated by
\begin{equation}\label{eq:Pabs-empirical}
  P_{\mathrm{abs}} \approx \Sinc \cdot T_{\mathrm{tr}}
  \cdot A_\perp(\khat),
\end{equation}
where $A_\perp$ is the projected (shadow) area of the body in the
direction of incidence $\khat$ and $T_{\mathrm{tr}}$ is a
transmission coefficient nearly independent of incidence angle and
polarisation. The empirical agreement is at the few-percent level.
Kodera et al.\ report $5\,\%$ across $10$ to $100\,$GHz against full
anatomical FDTD, and
Diao et al.\ report a $3\,\%$ discrepancy between FDTD and a Fresnel
estimate at $28\,$GHz. Yet $T_{\mathrm{tr}}$ is computed in every
case from a one-dimensional slab model fitted to FDTD output, and no
mechanism is given for its near-constancy across angle and
polarisation. The local map from incidence direction to absorbed
power on a curved body is not addressed, and the partial closed form
is treated as a numerical accident rather than a structural property
of biological tissue at mmWave.

This paper derives the local map from first principles. The result
is the closed-form local absorption law
\begin{equation}\label{eq:law}
  \Sab(\rr) = \Sinc \cdot T_0 \cdot \Vis(\rr,\khat) \cdot
  \pospart{\nhat(\rr) \cdot (-\khat)}.
\end{equation}
The positive-part operator $(\cdot)_+ \equiv \max(0, \cdot)$ gates the
geometric cosine. Front-facing surface elements absorb at the cosine
rate, and back-facing elements absorb nothing. The factor
$\Vis(\rr,\khat) \in \{0,1\}$ is the binary self-shadowing visibility,
equal to one when the surface point $\rr$ is reached by an
unobstructed ray from direction $\khat$ and zero when $\rr$ is
shadowed by the body itself. For a convex body $\Vis \equiv 1$ and
\eqref{eq:law} reduces to $\Sab(\rr) = \Sinc T_0 \pospart{\nhat\cdot
(-\khat)}$, which is the classical convex form. The factor $\Vis$ is
the ambient-occlusion factor of computer
graphics~\cite{Zhukov1998,Landis2002,AkenineMoller2018}, evaluated at
frame rate on consumer GPUs. In machine-learning notation, $(\cdot)_+$
is the ReLU activation; the diffraction correction in
\cref{sec:corrections} replaces it with a smoother analogue resembling
the GELU activation. The scalar $T_0$ depends only on the
tissue dielectric properties and is computed in closed form from the
complex refractive index. The transmission coefficient
$T_{\mathrm{tr}}$ used in the literature is identified with $T_0$. A second result names the
mechanism behind the near-constancy of $T_0$ across angle. A
polarisation-aware treatment of the surface flux gives an exact
identity in which the angular dependence of the unpolarised Fresnel
transmission $\Tavg(\theta)$ is the only obstruction to
\eqref{eq:law}. We show that $\Tavg(\theta)$ stays within $5.6\,\%$
of its normal-incidence value $T_0$ over $[0^\circ, 75^\circ]$ for
skin at $28\,$GHz, and within $5\,\%$ root-mean-square across $0.3$
to $100\,$GHz on all biological tissues with refractive index modulus
above $2.5$. The mechanism is a TM-rising and TE-falling compensation
that places lossy biological tissue in the high-index regime
described by Azzam~\cite{Azzam2015} for optical substrates. Its
application to dosimetry has not been made in either the optics or
the bioelectromagnetics literature. A third result is computational. For multiple incident waves with
directions $\khat_i$ and power densities $S_i$, the multi-source form
of \eqref{eq:law} is
\begin{equation}\label{eq:relu-net}
  \Sab(\rr) = T_0 \sum_{i=1}^N S_i\,\Vis(\rr,\khat_i)\,
  \pospart{\nhat(\rr) \cdot (-\khat_i)},
\end{equation}
the forward pass of a single-hidden-layer rectified network whose
weight matrix holds the negative incidence directions, whose output
weights are the tissue-scaled path powers, and whose hidden-unit gates
are the per-path visibilities, with an exact activation rather than a
fit. The chain from ray
tracer to dosimetry is differentiable end-to-end. The arithmetic primitive is
the same one a modern rendering engine runs at frame rate, and
gradients of regulatory quantities (whole-body specific absorption
rate $\mathrm{SAR}_{\mathrm{wb}}$, peak $\Sab$) with respect to
design parameters propagate by standard backpropagation.

The paper is organised as follows. \Cref{sec:law} derives the exact
polarisation-aware absorption law from Maxwell and identifies
$\Tavg(\theta)$ as the only remaining angle-dependent quantity.
\Cref{sec:pB} sets up the pseudo-Brewster compensation,
quantifies its accuracy across tissues and frequencies, and arrives
at the geometric law in \eqref{eq:law}. \Cref{sec:geo} states the
matrix form on a triangle mesh and the rectified-network identification.
\Cref{sec:val} compares the law against three independent
ground truths: exact Mie theory, full Fresnel integration on an
anatomical phantom, and Sim4Life FDTD on the same phantom.
\Cref{sec:corrections} treats higher-order corrections from
curvature, diffraction, and inter-body reflections.
\Cref{sec:disc} covers limitations and the regime of validity.


\section{Local absorption law}\label{sec:law}

The absorbed power density at a point on the body surface follows
from first principles for a plane wave on a locally flat tissue
interface. The derivation below is exact, and all approximations are
introduced in subsequent sections.

\subsection{Setup}
Consider a monochromatic plane wave with time-averaged Poynting
vector $\mathbf{S}_{\mathrm{inc}} = \Sinc \khat$ (units W/m$^2$)
incident on a body whose surface $\Sigma$ is locally flat on the
scale of the wavelength. The skin depth is much smaller than any
body dimension, so all power transmitted through the surface is
absorbed within a thin surface layer. We assume no coherent
reflections from internal tissue interfaces or from other body
parts. These assumptions are secondary at mmWave frequencies and are
discussed in \cref{sec:disc}.

At a surface point $\rr$ with outward unit normal $\nhat(\rr)$, the
incidence cosine is $\mu(\rr) \equiv \nhat(\rr) \cdot (-\khat) =
\cos\theta_i(\rr)$. A front-facing point has $\mu > 0$. A point
facing away from the source has $\mu \le 0$.

\subsection{Power flux through the surface}
The inward power flux through a surface element $\diff A$ at $\rr$
is $\diff P_{\mathrm{in}} = \Sinc \, \mu \, \diff A$. The reflected
wave propagates in the specular direction with power density
$|r(\theta_i)|^2 \Sinc$, where $r(\theta_i)$ is the Fresnel
amplitude reflection coefficient appropriate to the polarisation.
Energy conservation at a lossy half-space gives
\begin{equation}\label{eq:Sab-pol}
  \Sab(\rr) = \Sinc \cdot T(\theta_i) \cdot \mu(\rr),
  \qquad T(\theta_i) \equiv 1 - |r(\theta_i)|^2,
\end{equation}
for each polarisation separately. The quantity
$T(\theta_i)\,\mu(\rr)\,\Sinc$ is the transmitted power density
crossing the skin per unit skin area.

\subsection{Fresnel coefficients for a lossy half-space}
The body surface is modelled as a planar interface between free
space ($n_1 = 1$) and a lossy medium with complex refractive index
$\ntilde = \sqrt{\varepsilon_r - i\sigma/(\omega\varepsilon_0)}$.
The normal component of the wave vector inside the medium is
$\xi = \sqrt{\ntilde^2 - 1 + \mu^2}$ with $\RE(\xi) > 0$. Continuity
of the tangential electric and magnetic fields at the interface
gives the standard Fresnel coefficients
\begin{equation}\label{eq:rs-rp}
  r_s = \frac{\mu - \xi}{\mu + \xi},
  \qquad
  r_p = \frac{\ntilde^2\,\mu - \xi}{\ntilde^2\,\mu + \xi},
\end{equation}
for transverse-electric (TE) and transverse-magnetic (TM)
polarisations, respectively. The corresponding power-absorption
coefficients are $T_s(\theta) = 1 - |r_s|^2$ and $T_p(\theta) = 1 -
|r_p|^2$. At normal incidence, $\mu = 1$ and $\xi = \ntilde$, giving
the polarisation-degenerate value
\begin{equation}\label{eq:T0}
  T_0 \equiv T_s(0) = T_p(0) = \frac{4\,\RE(\ntilde)}{|1+\ntilde|^2}.
\end{equation}
For skin at $28\,$GHz with $\varepsilon_r = 16.55$ and
$\sigma = 25.8\,$S/m, $\ntilde = 4.49 - 1.79i$ and $T_0 = 0.536$.

\subsection{Polarisation-aware exact law}\label{subsec:exact-law}
A monochromatic plane wave is fully polarised. At a surface point
$\rr$, decompose the incident electric field into local TE and TM
components by projecting on the unit vectors
$\hat{e}_s(\rr) = \khat \times \nhat / |\khat \times \nhat|$ and
$\hat{e}_p(\rr) = \hat{e}_s \times \khat$. Writing the field as
$\EE_0 = a_s \hat{e}_s + a_p \hat{e}_p$ and the local TE and TM
energy fractions as $|e_s|^2 = |a_s|^2 / |\EE_0|^2$ and
$|e_p|^2 = |a_p|^2 / |\EE_0|^2$, the effective transmission at $\rr$
is
\begin{equation}\label{eq:Teff}
  \Teff(\rr) = |e_s(\rr)|^2 \, T_s(\theta) + |e_p(\rr)|^2 \,
  T_p(\theta).
\end{equation}
The exact absorbed power density at a visible point is therefore
\begin{equation}\label{eq:Sab-exact}
  \boxed{
  \Sab(\rr) = \Sinc \, \Teff(\rr) \, \pospart{\mu(\rr)},
  }
\end{equation}
where $(\cdot)_+$ gates back-facing elements to zero. \Cref{eq:Sab-exact}
holds for any polarisation, any frequency where the body is opaque,
and any locally flat surface. The self-shadowing factor $\Vis(\rr,\khat)$
of \eqref{eq:law} is suppressed in this and the next subsection because
the Fresnel calculation operates at a point already taken to be
visible. Visibility re-enters with the multi-source matrix form in
\cref{sec:geo}.

To proceed, write
\begin{equation}\label{eq:Teff-decomp}
  \Teff(\rr) = \Tavg(\theta) + \tfrac{1}{2}\,q(\rr)\,\Delta T(\theta),
\end{equation}
with $\Tavg = \tfrac{1}{2}(T_s + T_p)$ the unpolarised baseline,
$\Delta T = T_p - T_s$ the polarisation splitting, and $q = |e_p|^2
- |e_s|^2 \in [-1, 1]$ the local TM excess. Three conditions each
kill the second term. For circularly polarised illumination the
field's TE and TM intensities are equal at every point on every
body, and $q \equiv 0$ pointwise. For a monochromatic linearly
polarised plane wave with random ensemble orientation the average
$\langle q \rangle$ vanishes by the same argument. For a body in a
richly multipath environment with $N$ independent path orientations,
the variance of the polarisation correction scales as
$D_B / \sqrt{2N}$, where $D_B$ is the body's polarisation directivity.
For $N \ge 20$ paths the correction drops below $2.5\,\%$. Under any
of these conditions the exact law reduces to
\begin{equation}\label{eq:Sab-Tavg}
  \Sab(\rr) = \Sinc \, \Tavg(\theta(\rr)) \, \pospart{\mu(\rr)}.
\end{equation}
The angular dependence is now confined to the scalar function
$\Tavg(\theta)$. The next section shows that this function is nearly
constant for biological tissue.


\section{Pseudo-Brewster compensation}\label{sec:pB}

The unpolarised Fresnel transmission $\Tavg(\theta)$ is the average
of two angularly opposed quantities. As the incidence angle
increases, $T_s(\theta)$ decreases monotonically, while $T_p(\theta)$
rises towards a near-unity peak at the pseudo-Brewster angle and
then falls. Their average tracks $T_0$ across a wide angular range.
For biological tissue at mmWave frequencies, the angle-dependent
factor in \eqref{eq:Sab-Tavg} reduces to the scalar $T_0$ at the
cost of an error of a few percent, leaving the absorbed-power
density as a function of body geometry alone.

\subsection{Physical mechanism}\label{subsec:pB-mech}
The Brewster angle of a lossless dielectric is $\theta_B =
\arctan(n_2/n_1)$, at which the TM reflection coefficient vanishes
\cite{BornWolf1999}. For a lossy dielectric the reflection minimum
is finite but small. The angle at which $|r_p|^2$ is minimised is
called the pseudo-Brewster angle and satisfies $\theta_{pB} \approx
\arctan|\ntilde|$ to within $1^\circ$ for $|\ntilde| > 3$
\cite{Potter1970,Ohman1977}. At this angle, $T_p$ reaches a peak
that is often above $0.95$, while $T_s$ has fallen below $0.20$.
Their average $\Tavg(\theta_{pB}) \approx 0.5$ is close to the
normal-incidence value $T_0 \approx 0.5$ to $0.6$ for biological
tissue at mmWave.

Azzam \cite{Azzam2015} showed that for lossless dielectric
substrates with refractive index $|\ntilde| > 2 + \sqrt{3} \approx
3.73$, the unpolarised reflectance varies by less than $1\,\%$ over
$[0^\circ, 60^\circ]$. The criterion weakens for $|\ntilde| > 2.5$
but the near-constancy persists. Biological tissue at mmWave
frequencies has $|\ntilde| \in [3, 6]$, placing it in this regime.
The application of the high-index criterion to biological dosimetry
has not been made in the literature.

\subsection{Quantitative behaviour across angle and polarisation}
\label{subsec:pB-quant}
\Cref{fig:apd-angle} illustrates the compensation for skin at
$28\,$GHz. The upper panel shows $T_s$, $T_p$, and $\Tavg$ as
functions of incidence angle. The lower panel shows the absorbed
power density $\Sab/\Sinc = T(\theta)\cos\theta$ for each
polarisation and for the simplified product $T_0 \cos\theta$. The
unpolarised curve closely tracks the simplified prediction, and the
small gap is the Fresnel approximation error.

\begin{figure}[!t]
  \centering
  \includegraphics[width=\columnwidth]{apd_angle_dependence.pdf}
  \caption{Pseudo-Brewster compensation for skin at $28\,$GHz
  ($\ntilde = 4.49 - 1.79i$, $T_0 = 0.539$). Upper panel: Fresnel
  power-absorption coefficients $T_s$ (TE), $T_p$ (TM), and
  $\Tavg = \tfrac{1}{2}(T_s + T_p)$ versus incidence angle. The
  TE and TM curves diverge strongly, and their average stays within
  $5.6\,\%$ of $T_0$ up to $75^\circ$. Lower panel: normalised
  absorbed power density $\Sab/\Sinc = T(\theta)\cos\theta$. The
  dashed grey curve is the simplified $T_0 \cos\theta$ prediction.
  The gap between the dashed grey curve and the green unpolarised
  curve is the entire Fresnel approximation error.}
  \label{fig:apd-angle}
\end{figure}

\Cref{tab:fresnel-skin} quantifies the deviation of $\Tavg$ from
$T_0$ across $[0^\circ, 75^\circ]$. The maximum is $5.6\,\%$ at $70$
to $75^\circ$. Below $30^\circ$ the agreement is at the fourth
significant figure.

\begin{table}[!t]
\centering
\caption{Fresnel transmission for skin at $28\,$GHz. The ratio
$\Tavg/T_0$ stays within $5.6\,\%$ of unity over $[0^\circ,
75^\circ]$.}
\label{tab:fresnel-skin}
\begin{tabular}{ccccc}
\toprule
$\theta$ & $T_s$ (TE) & $T_p$ (TM) & $\Tavg$ & $\Tavg/T_0$ \\
\midrule
$0^\circ$  & 0.539 & 0.539 & 0.539 & 1.000 \\
$30^\circ$ & 0.489 & 0.591 & 0.540 & 1.002 \\
$45^\circ$ & 0.422 & 0.666 & 0.544 & 1.010 \\
$60^\circ$ & 0.321 & 0.791 & 0.556 & 1.032 \\
$70^\circ$ & 0.233 & 0.902 & 0.568 & 1.054 \\
$75^\circ$ & 0.182 & 0.952 & 0.567 & 1.053 \\
\bottomrule
\end{tabular}
\end{table}

\subsection{Tissue universality}\label{subsec:pB-tissues}
All biological tissues at mmWave frequencies cluster in the
$|\ntilde| > 2.5$ region where the compensation operates.
\Cref{tab:materials} lists the relevant parameters. Skin and muscle
have $|\ntilde|$ near $5$ and a $\Tavg/T_0$ variation below
$5.6\,\%$. Water has $|\ntilde| > 6$ and a variation below $4\,\%$.
Fat is the outlier, with $|\ntilde| \approx 2$ and an $8.2\,\%$
variation, but fat is rarely the outermost tissue layer at exposure
sites of regulatory interest.

\begin{table}[!t]
\centering
\caption{Pseudo-Brewster compensation across tissue types at
$28\,$GHz. The variation column is the maximum deviation of
$\Tavg/T_0$ from unity over $[0^\circ, 75^\circ]$.}
\label{tab:materials}
\begin{tabular}{lccccc}
\toprule
Tissue & $\varepsilon_r$ & $\sigma$ (S/m) & $|\ntilde|$ & $T_0$
  & $\Tavg$ var. \\
\midrule
Skin   & 17.0 & 25.0 & 4.84 & 0.54 & $5.6\,\%$ \\
Muscle & 25.0 & 30.0 & 5.62 & 0.48 & $4.8\,\%$ \\
Fat    &  4.0 &  2.0 & 2.05 & 0.77 & $8.2\,\%$ \\
Water  & 25.0 & 55.0 & 6.62 & 0.45 & $3.9\,\%$ \\
\bottomrule
\end{tabular}
\end{table}

\subsection{Frequency landscape}\label{subsec:pB-freq}
The accuracy of the constant-$T_0$ approximation has a clean
frequency dependence. Define the sphere ratio
\begin{equation}\label{eq:Tbar-def}
  R(f) \equiv \frac{T_0(f)}{\Tbar(f)},
  \qquad
  \Tbar(f) \equiv 2 \int_0^1 \Tavg(\mu, f)\,\mu\,\diff\mu,
\end{equation}
where $\Tbar$ is the flux-weighted average transmission. $R(f) = 1$
when the constant-$T_0$ approximation is exact on a
direction-averaged quantity. $R < 1$ means $T_0$ underestimates the
absorbed power, and $R > 1$ means it overestimates. \Cref{fig:R-of-f}
shows that for skin the crossover occurs at $40.4\,$GHz, where the
compensation is exact to four significant figures. Below this
frequency the approximation is conservative (it underestimates
absorbed power, which is favourable for compliance). Above, it
overestimates by at most $3.5\,\%$ at $100\,$GHz. The
root-mean-square error remains below $5\,\%$ across $0.3$ to
$100\,$GHz.

\begin{figure}[!t]
  \centering
  \includegraphics[width=\columnwidth]{R_of_f_landscape.pdf}
  \caption{Pseudo-Brewster compensation across frequency for skin
  (IT'IS v5.0 dielectric model). Top panel: ratio
  $R(f) = T_0/\Tbar$. The compensation is exact at $40.4\,$GHz, and
  the conservative region ($R < 1$, $T_0$ underestimates) extends
  below this frequency. Middle panel: $T_0(f)$ and $\Tbar(f)$.
  Bottom panel: activation function $\Tavg(\theta)\cos\theta$ at six
  frequencies. The dashed line is the $T_0 \cos\theta$ approximation
  evaluated at $40\,$GHz, where the curves nearly overlap.}
  \label{fig:R-of-f}
\end{figure}

\subsection{Geometric absorption law}\label{subsec:pB-geom}
Combining the polarisation reduction \eqref{eq:Sab-Tavg} with the
pseudo-Brewster substitution $\Tavg(\theta) \to T_0$ and reinstating
self-shadowing through the binary visibility $\Vis(\rr,\khat)$, the
exact absorption law \eqref{eq:Sab-exact} reduces to the geometric
form
\begin{equation}\label{eq:geom-law}
  \boxed{
  \Sab(\rr) \approx \Sinc \cdot T_0 \cdot \Vis(\rr,\khat) \cdot
  \pospart{\nhat(\rr) \cdot (-\khat)}.
  }
\end{equation}
The tissue physics enters through the scalar $T_0$. All spatial
variation is determined by the body's shape through the
surface-normal field $\nhat(\rr)$ and the visibility field
$\Vis(\rr,\khat)$. For a convex body $\Vis \equiv 1$ and
\eqref{eq:geom-law} reduces to the classical convex form
$\Sinc T_0 \pospart{\nhat\cdot(-\khat)}$. \Cref{fig:sab-phantom}
contrasts the convex-limit prediction with the visibility-aware
result on the Thelonious phantom and motivates the explicit
$\Vis$ factor. The transmission coefficient
$T_{\mathrm{tr}}$ fitted in
\cite{Kodera2024,Diao2024,Funahashi2018} is identified with $T_0$,
and its near-constancy across angle and polarisation is identified
with the pseudo-Brewster compensation.

\begin{figure*}[!t]
  \centering
  \includegraphics[width=\textwidth]{sab_phantom_pair.pdf}
  \caption{Absorbed power density $\Sab(\rr)$ on the Thelonious
  anatomical phantom at $28\,$GHz, $\Sinc = 1\,$W/m$^2$, frontal
  illumination $\khat = +\hat{y}$. Panel (a) is the convex-limit
  prediction with $\Vis \equiv 1$: every front-facing element
  absorbs at the cosine rate, including elements that the body
  itself shadows (the underside of the chin, the medial sides of the
  thighs, the inside of the wrists). Panel (b) is the
  visibility-aware result with $\Vis(\rr,\khat)$ from the
  ambient-occlusion solver: shadowed triangles drop to zero. The
  two panels share the colour scale $[0, T_0\,\Sinc]$. The contrast
  motivates the explicit visibility factor in \eqref{eq:geom-law};
  the convex form is a useful reference but is not faithful on a
  non-convex body.}
  \label{fig:sab-phantom}
\end{figure*}

The error of \eqref{eq:geom-law} relative to the exact
polarisation-aware law is bounded by the maximum of the polarisation
correction and the angular variation of $\Tavg$. The latter is below
$5\,\%$ across the mmWave band. The former is below $2.5\,\%$ for
$N \ge 20$ independent multipath components, and at most
$16\,\%$ on a real human body for the worst-case linearly polarised
single plane wave. The combined error is below the $20\,\%$
uncertainty in tissue dielectric properties at mmWave
\cite{AlekseevZiskin2007}.


\section{Geometric reduction and computational form}\label{sec:geo}

The geometric law \eqref{eq:geom-law} reduces dosimetry to a
function of body geometry alone. This section gives the matrix form
on a triangle mesh, the multi-source extension, and the rectified-network
identification that makes the dosimetry chain differentiable.

\subsection{Matrix form on a triangle mesh}
Discretise the body surface as a triangle mesh with $M$ triangles.
Let $\mathbf{N} \in \mathbb{R}^{M \times 3}$ be the matrix of
outward unit normals (one row per triangle). For multiple incident
plane waves with directions $\khat_i \in \mathbb{R}^3$ and power
densities $S_i \in \mathbb{R}_+$, $i = 1, \ldots, N$, define the
direction matrix $\mathbf{K} = [-\khat_1, \ldots, -\khat_N] \in
\mathbb{R}^{3 \times N}$ and the power vector $\mathbf{s} = [S_1,
\ldots, S_N]^\top$. Let $\mathbf{V} \in \{0, 1\}^{M \times N}$ be the
visibility matrix with entries $V_{ji} = \Vis(\rr_j, \khat_i)$ from
ray casting against the mesh, $\rr_j$ the centroid of triangle $j$.
The per-triangle absorbed power density vector is
\begin{equation}\label{eq:mat-multi}
  \mathbf{S}_{\mathrm{ab}} = T_0\,\bigl(\pospart{\mathbf{N}\,\mathbf{K}}
  \odot \mathbf{V}\bigr)\,\mathbf{s},
\end{equation}
where $(\cdot)_+$ acts componentwise on the cosine matrix
$\mathbf{N}\,\mathbf{K} \in \mathbb{R}^{M \times N}$, gating
back-facing combinations, and the Hadamard product with $\mathbf{V}$
gates self-shadowed combinations. For a single incidence
direction $\khat$,
\begin{equation}\label{eq:mat-single}
  \mathbf{S}_{\mathrm{ab}} = \Sinc \, T_0 \,
  \bigl(\pospart{\mathbf{N}\,(-\khat)} \odot \mathbf{v}(\khat)\bigr),
\end{equation}
with $\mathbf{v}(\khat) \in \{0,1\}^M$ the column of $\mathbf{V}$ for
direction $\khat$. For convex bodies $\mathbf{V} = \mathbf{1}$ and
the Hadamard product drops out.

The ICNIRP $4\,\mathrm{cm}^2$ spatial averaging is a sparse matrix
$\mathbf{G} \in \mathbb{R}^{M \times M}$ with row-stochastic
neighbourhood structure. The averaged map is
$\bar{\mathbf{S}}_{\mathrm{ab}} = \mathbf{G}\,\mathbf{S}_{\mathrm{ab}}$.
Total absorbed power follows from $P_{\mathrm{abs}} = \mathbf{a}^\top
\mathbf{S}_{\mathrm{ab}}$, where $\mathbf{a}$ is the triangle-area
vector. Whole-body SAR is $\mathrm{SAR}_{\mathrm{wb}} =
P_{\mathrm{abs}} / m$ for body mass $m$.

\subsection{Hierarchical computation}
The full per-triangle map costs $O(MN)$. When only aggregate
quantities are needed, the projected-area identity $P_{\mathrm{abs}}
= \Sinc\,T_0\,A_\perp(\khat)$ is much faster. Define the
projected (shadow) area $A_\perp(\khat) = \int_{\Sigma_+}
\cos\theta\,\diff A$ as the area of the body's silhouette
perpendicular to $\khat$. For multiple sources,
\begin{equation}\label{eq:Pabs-aggregate}
  P_{\mathrm{abs}} = T_0 \sum_{i=1}^N S_i\,A_\perp(\khat_i),
\end{equation}
which is $O(N)$ when $A_\perp$ is precomputed as a lookup table on
the unit sphere. The body mesh does not enter the inner loop. A
spherical-harmonic representation of $A_\perp(\khat)$ at degree $L
\le 6$ achieves sub-percent accuracy and reduces the lookup to a dot
product of coefficient vectors.

The two regimes give a natural hierarchy. The bound regime ($\Sinc
\le \Sinc^{\max}$) costs $O(1)$ per scenario and uses the worst-case
projected-area bound $A_\perp \le A/2$. The aggregate regime
\eqref{eq:Pabs-aggregate} gives exact total power and SAR for $O(N)$
per scenario. The spatial regime \eqref{eq:mat-multi} gives the full
per-triangle map for $O(MN)$. With $M \sim 10^4$ and $N \sim 10^2$,
the aggregate regime is roughly $10^4$ times faster than the spatial
regime.

\subsection{Rectified network interpretation}
\Cref{eq:mat-multi} is the forward pass of a single-hidden-layer
network. The hidden layer has $N$ units. The weight matrix
$\mathbf{W} = \mathbf{K}^\top \in \mathbb{R}^{N \times 3}$ is the
matrix of negative incidence directions. The output weights are the
tissue-scaled path powers $T_0\,\mathbf{s} \in \mathbb{R}^N$. The
bias is zero. The activation is the positive part $(\cdot)_+$, which
in machine-learning notation is the ReLU activation, applied here
with an additional binary visibility mask $\mathbf{V}$ that gates
self-shadowed inputs. The network evaluates at every surface point
$\rr$ by taking $\nhat(\rr)$ as input. The activation is exact rather
than fitted, and the structure follows from the geometric law rather
than from a training procedure. Three consequences follow.

First, the chain from ray tracer to dosimetry is differentiable
end-to-end. Each block (path tracing, matrix multiplication,
positive-part gate, weighted sum) is differentiable. Gradients of regulatory quantities
with respect to design parameters (antenna positions, antenna
orientations, RIS phases) propagate by standard backpropagation.
Modern GPU-accelerated ray-tracing engines such as Sionna~RT
\cite{SionnaRT} give differentiable propagation, and the
dosimetry layer attaches as a single matrix-multiply-and-clip pass.

Second, the arithmetic primitive is the same one a rendering engine
runs at frame rate. The Hadamard product
$\pospart{\mathbf{N}\,\mathbf{K}} \odot \mathbf{V}$ is a per-pixel
shading operation: the cosine gate is rectified shading and
$\mathbf{V}$ is ambient occlusion. For a body mesh of $M \sim 10^4$ triangles and
$N \sim 10^2$ paths, the spatial map is one matrix-vector multiply
on the GPU, so real-time exposure monitoring runs on
consumer hardware.

Third, gradient-based optimisation of antenna configurations under
exposure constraints is a backpropagation problem in the space of
path powers and directions. Problems that previously
needed gradient-free methods (genetic algorithms, particle
swarm) reduce to standard differentiable optimisation in a
low-dimensional ambient space.


\section{Validation}\label{sec:val}

We test the law against three independent ground truths.
The first is exact Mie theory for lossy dielectric spheres, which
tests it on a smooth shape with no curvature variation
across the surface. The second is full Fresnel integration on the
Thelonious anatomical phantom, which tests it on
realistic body geometry. The third is full Sim4Life FDTD on the same
phantom with matched dielectric properties and matched plane-wave
excitation. The three together span the regimes (analytic,
semi-analytic, volumetric numerical) and the geometries (sphere,
anatomy, both) that bound the claimed window of
validity.

\subsection{Mie theory on lossy spheres}
For a lossy sphere of radius $a$ and complex refractive index
$\ntilde$, the Mie series gives an exact solution for the absorption
efficiency $Q_{\mathrm{abs}}$. The law predicts
$P_{\mathrm{abs}} = \Sinc \, T_0 \, \pi a^2$, so its error
is $(T_0 / Q_{\mathrm{abs}} - 1) \times 100\,\%$. We use IT'IS v5.0
skin properties at each frequency. The total error
splits into two contributions. The Fresnel approximation error
is shape- and frequency-dependent but size-independent. On
a sphere, it is the sphere ratio $R = T_0 / \langle
\Tavg \rangle$, which crosses unity at $39\,$GHz (\cref{fig:R-of-f}).
The diffraction error is size-dependent and scales as $x^{-2/3}$ in
the optical regime, where $x = \pi d / \lambda$ is the size
parameter. Diffraction bends waves into the geometric shadow,
adding absorption that the law misses.

\Cref{fig:mie} collects the Mie validation. The upper panel shows
the error versus size parameter at $28\,$GHz. It
converges from below towards the Fresnel limit
$R_{\mathrm{sphere}} - 1 \approx -1.2\,\%$ as $x \to \infty$. The
middle panel shows the error versus frequency for four
representative body-part diameters. The lower panel shows the
frequency-dependent Fresnel limit $R_{\mathrm{sphere}}(f)$.

\begin{figure}[!t]
  \centering
  \includegraphics[width=\columnwidth]{mie_validation_corrected.pdf}
  \caption{Mie validation against lossy spheres with
  frequency-dependent IT'IS skin properties. Upper panel: prediction
  error versus size parameter at $28\,$GHz. Vertical dashed lines
  mark body-part sizes. Middle panel: error versus
  frequency for finger ($17\,$mm), arm ($80\,$mm), head ($180\,$mm),
  and torso ($300\,$mm) diameters. The mmWave band is shaded green.
  Lower panel: frequency-dependent Fresnel limit
  $R_{\mathrm{sphere}}(f) = T_0/\langle\Tavg\rangle$, crossing unity
  at $\sim 39\,$GHz.}
  \label{fig:mie}
\end{figure}

For body-relevant sizes (head, torso) at mmWave, the error
is $3$ to $14\,\%$, set mostly by diffraction into the geometric
shadow, and the law underestimates absorption (a conservative
error in the regulatory sense). Fingers at sub-mmWave frequencies
are the worst case, with errors above $30\,\%$. The error stays
below the $20\,\%$ tissue dielectric uncertainty for body-scale
objects in the claimed regime.

\subsection{Full Fresnel on the Thelonious phantom}
The Mie test bounds the Fresnel error on a smooth shape. The
phantom test bounds it on realistic anatomy. We compare the
simplified prediction $\Sab^{\mathrm{simp}} = \Sinc T_0 \pospart{\mu}$
against the full polarisation-aware Fresnel integration
$\Sab^{\mathrm{full}} = \Sinc \Teff(\theta, \mathrm{pol})
\pospart{\mu}$ on the Thelonious mesh ($23\,826$ triangles,
$0.787\,\mathrm{m}^2$ surface area). The incident plane wave comes
from above, with skin properties at $28\,$GHz.

\begin{table}[!t]
\centering
\caption{Geometric law versus full Fresnel integration on the Thelonious
phantom (skin at $28\,$GHz, plane wave from above).}
\label{tab:phantom}
\begin{tabular}{lccc}
\toprule
Metric & Simplified & Full Fresnel & Error \\
\midrule
Total absorbed power     & $48.12\,$mW    & $47.95\,$mW    & $0.35\,\%$ \\
Mean $\Sab$ (illum.)     & $0.185\,$W/m$^2$ & $0.186\,$W/m$^2$ & $0.5\,\%$ \\
Peak $\Sab$              & $0.539\,$W/m$^2$ & $0.539\,$W/m$^2$ & $0.0\,\%$ \\
\bottomrule
\end{tabular}
\end{table}

For the $4\,907$ illuminated triangles with $\theta < 75^\circ$, the
local statistics are mean error $-2.6\,\%$, RMS $3.2\,\%$, range
$[-5.3\,\%, 0.0\,\%]$. The peak $\Sab$ is recovered exactly because
the maximum of $\Sab$ occurs at normal incidence, where $\Teff(0) =
T_0$ regardless of polarisation. The total power error is below
$1\,\%$. The local error is below $5.5\,\%$ everywhere with $\theta
< 75^\circ$. \Cref{fig:apd-direction} extends the analysis to $128$
illumination directions and three polarisation states.

\begin{figure*}[!t]
  \centering
  \includegraphics[width=\textwidth]{apd_direction_analysis.pdf}
  \caption{Directional and polarisation analysis on the Thelonious
  phantom (skin at $28\,$GHz, $\Sinc = 1\,$W/m$^2$). Upper-left:
  total absorbed power across $128$ directions for TE (red), TM
  (blue), and unpolarised (green). Upper-right: probability density
  of local $\Sab$ on illuminated triangles. Lower-left: box plots by
  polarisation, with the dashed line marking the prediction
  $T_0 A_\perp = 105.9\,$mW. Lower-right: coefficient of variation
  of total absorbed power and peak $\Sab$ across directions.}
  \label{fig:apd-direction}
\end{figure*}

\subsection{Full FDTD on the Thelonious phantom}
The Mie and phantom tests check the Fresnel approximation against
analytic and semi-analytic ground truths. We complete the validation
with full Sim4Life FDTD on the same Thelonious mesh, matched
dielectric properties, and matched plane-wave excitation.

The IEC/IEEE\,$63195$ regulatory metric is the peak $\Sab$ averaged
over a $4\,\mathrm{cm}^2$ patch on the body surface. We compare the
prediction of this peak against Sim4Life's
\texttt{GenericSAPDEvaluator} on the same scenarios. At $7\,$GHz on
the Thelonious phantom, with three lateral and frontal incidence
directions and $\theta$-polarisation, the direction-averaged ratio
of law to FDTD on the peak $4\,\mathrm{cm}^2$ metric is
$1.027$. Per-direction values are $1.06$, $1.20$, and $0.83$. All
three lie within the $[0.80, 1.20]$ acceptance band that the
regulatory literature uses to certify computational electromagnetic
models for mmWave dosimetry.

The total absorbed power on the body integrated from FDTD's
volumetric $\sigma|E|^2/2$ output gives a different ratio of $0.479$
on the same three scenarios. The ratio at the closely related
$5.8\,$GHz endpoint is consistent with $1.0$ to within a few
percent, suggesting that the $7\,$GHz total-power discrepancy comes
from FDTD-side post-processing rather than from a breakdown of the
law. We retain only the peak $4\,\mathrm{cm}^2$
metric as the validated comparison and note the open question for
future work.

The third numerical check is the direction-averaged Cauchy formula
$\langle P_{\mathrm{abs}} \rangle = \Sinc \Tbar A_{\mathrm{ab}}/4$
applied to predictions across $12$ directions and $2$ polarisations
at $5.8\,$GHz. The ratio of law to FDTD on direction-averaged total
absorbed power is $1.012$. The
non-convexity correction $A_{\mathrm{ab}}/A = 0.865$ used in this
calculation comes from an ambient-occlusion solver on the same mesh
and matches the body's classical $A/4$ silhouette ratio to within
$0.6\,\%$ \cite{Cauchy1841}. This is an internal consistency check
on the geometric pipeline.

\Cref{fig:val-fdtd} summarises the FDTD comparison across the
sub-$6\,$GHz band. The closed-form Cauchy prediction (black stars)
approaches unity at the upper end of the band, where the geometric
optics asymptotic is expected to apply. Below $5\,$GHz the
law underestimates because body-scale Mie and resonance
effects do not enter a surface-only law. The breakdown is
consistent with the Mie analysis on a sphere of comparable size
parameter.

\begin{figure*}[!t]
  \centering
  \includegraphics[width=\textwidth]{fig_kernels_vs_fdtd.png}
  \caption{Direction-averaged comparison of the geometric law against
  full Sim4Life FDTD on the Thelonious phantom across $0.45$ to
  $5.8\,$GHz, $12$ directions and $2$ polarisations per frequency.
  The legend turns on the corrections of \cref{sec:corrections} in
  sequence: ``Fresnel only'' uses the angle-dependent Fresnel
  $T_s/T_p$ at the convex limit, ``+ polarisation'' adds the
  polarisation-aware $\Teff$, ``+ curvature \& diffraction'' adds
  the $1/(kR)$ refinement, ``Full kernel'' uses the constant-$T_0$
  geometric law without occlusion, ``+ occlusion'' multiplies by
  $\Vis(\rr,\khat)$, and ``Full + occlusion'' is the constant-$T_0$
  visibility-gated form. The closed-form Cauchy prediction (black
  stars, right panel) gives $\langle P_{\mathrm{abs}}\rangle^{\mathrm{law}}
  / \langle P_{\mathrm{abs}}\rangle^{\mathrm{FDTD}} = 1.012$ at
  $5.8\,$GHz with no fitted parameters. Below $5\,$GHz the
  law underestimates because body-scale Mie and resonance
  effects do not enter a surface-only law.}
  \label{fig:val-fdtd}
\end{figure*}


\section{Higher-order corrections}\label{sec:corrections}

The geometric law \eqref{eq:geom-law} treats the surface as locally
flat, gates back-facing elements with a sharp $(\cdot)_+$ cutoff, and
ignores power that re-enters the body after a first reflection. Three
systematic corrections refine each of these. The labels in
\cref{fig:val-fdtd} (``Fresnel only'', ``+ polarisation'',
``+ curvature \& diffraction'', ``Full kernel'', ``+ occlusion'') turn
each correction on in turn against the same FDTD reference.

\subsection{Curvature}\label{subsec:corr-curv}

For a surface with twice the local mean curvature $H = 1/R_1 + 1/R_2$,
the first-order Physical Optics correction multiplies the geometric
law by $1 + \mu/(kR_1) + \mu/(kR_2)$, where $k = 2\pi/\lambda$ is the
free-space wavenumber. Since $\pospart{\mu}\cdot\mu = \pospart{\mu}^2$,
the per-triangle update separates additively,
\begin{equation}\label{eq:curv-update}
  \Sab^{(j)} = T_0 \sum_i S_i\!\left[
  \pospart{\mu_{ji}} + \frac{H_j}{k}\,\pospart{\mu_{ji}}^2 \right]
  V_{ji},
\end{equation}
adding a quadratic gate on top of the linear one. The magnitude is set
by $1/(kR)$. \Cref{tab:curv-mag} lists the correction at $28\,$GHz on
representative body parts.

\begin{table}[!t]
\centering
\caption{Curvature correction at $28\,$GHz ($k \approx 587\,$m$^{-1}$).
The correction is below the Fresnel approximation error for most
body regions and concentrates at small features.}
\label{tab:curv-mag}
\begin{tabular}{lccc}
\toprule
Region & $R$\,(mm) & $1/(kR)$ & Correction \\
\midrule
Torso, head & $> 100$ & $< 0.2\,\%$ & negligible \\
Arm         & $\sim 40$ & $0.4\,\%$ & $\sim 0.4\,\%$ \\
Finger      & $\sim 8$  & $2.1\,\%$ & $\sim 2\,\%$ \\
Ear edge    & $\sim 2$  & $8.5\,\%$ & $\sim 8\,\%$ \\
\bottomrule
\end{tabular}
\end{table}

The correction grows as the wavelength approaches the local
body-part size. At sub-$6\,$GHz frequencies the smallest features
(fingers, ears) sit in $kR \lesssim 5$ where the correction is no
longer small. At $28\,$GHz, only the ear edges and finger tips carry
a correction above the Fresnel error floor. The middle curve in
\cref{fig:val-fdtd} (``+ curvature \& diffraction'') closes the gap
to the FDTD reference at the upper end of the band.

\subsection{Diffraction at the shadow boundary}\label{subsec:corr-diff}

The sharp $(\cdot)_+$ cutoff at $\mu = 0$ is a geometric-optics
idealisation. Diffraction smooths the shadow edge over a Fresnel-zone
width $\sqrt{\lambda R_j}$. The physical activation becomes
\begin{equation}\label{eq:gelu}
  \mu \mapsto \mu\;\tfrac{1}{2}\!\left[1 + \mathrm{erf}(\mu/\sigma_j)\right],
  \qquad \sigma_j = \sqrt{\lambda / (2\pi R_j)},
\end{equation}
which in machine-learning notation is the GELU activation, with the
smoothing scale set by physics rather than learned. The ML
analogy of $(\cdot)_+ \to$ ReLU now becomes literal: the geometric
law is the ReLU limit, and diffraction replaces it with GELU near
the shadow boundary. ICNIRP's $4\,\mathrm{cm}^2$ spatial averaging
already regularises the boundary on the scale that the regulator
cares about, so the GELU smoothing matters only for high-resolution
local maps or comparisons against point measurements. The integrated
effect on whole-body absorbed power is sub-$1\,\%$ above $10\,$GHz on
the Thelonious phantom and grows to a few percent below $6\,$GHz, in
line with the Mie analysis on body-scale spheres in \cref{sec:val}.

\subsection{Inter-body reflections}\label{subsec:corr-interbody}

After the body absorbs at one point, the reflected fraction
$1 - T_0 \approx 0.46$ propagates away. On a non-convex body, a part
of this reflected power re-illuminates another point and contributes
to absorption that the first-bounce law misses. The radiosity series
gives a multiplier $C(\rr) = 1/(1 - \bar{R}\,f(\rr))$ at each point,
where $\bar{R} = 1 - \Tbar \approx 0.46$ is the flux-weighted
reflectance and $f(\rr) \le 1 - \eta(\rr)$ is the recapture fraction
bounded by the local non-visible hemisphere area. Two effects keep
the body-averaged correction small. The bound $f \le 1 - \eta$
self-compensates: deep concavities ($\eta$ low) recapture more
($f$ high), so the product $\eta \cdot C$ moves much less than
$\eta$ alone. Specular reflection at mmWave (skin meets the Rayleigh
roughness criterion) reduces $f$ by roughly a factor of three
relative to the diffuse bound. On the Thelonious phantom at
$28\,$GHz, the area-weighted recapture fraction is
$f_{\mathrm{global}} \approx 0.09$ under the diffuse bound, giving
$C \approx 1.04$; the specular estimate brings this to
$C \approx 1.01$. The body-averaged correction stays below $2\,\%$,
small against the dielectric uncertainty of $20\,\%$ that dominates
the error budget. Correspondingly, the convex-hull energy bound
$\langle P_{\mathrm{abs}} \rangle \le \Sinc A_{\mathrm{CH}}/4$
brackets the true absorbed power within $A_{\mathrm{CH}}/A \approx
1.20$ on Thelonious, well within the dielectric floor.

\subsection{Where each correction matters}\label{subsec:corr-summary}

\Cref{fig:err-budget} collects the error contributions. The Fresnel
approximation that this paper bounds is below $5.6\,\%$ on skin at
$28\,$GHz. Curvature and diffraction together account for the
sub-$6\,$GHz residual visible in \cref{fig:val-fdtd}. Inter-body
reflections sit below $2\,\%$ at body-averaged level. The tissue
dielectric data is the dominant uncertainty at $\sim 20\,\%$
\cite{AlekseevZiskin2007} and propagates linearly into $T_0$.

\begin{figure}[!t]
  \centering
  \includegraphics[width=\columnwidth]{error_budget_comprehensive.pdf}
  \caption{Error budget for the geometric absorption law on body-scale
  geometries at mmWave. The Fresnel approximation error (this paper)
  is bounded by $5.6\,\%$ on skin at $28\,$GHz. Diffraction at body
  edges contributes up to $\sim 14\,\%$ at body-scale features. The
  tissue dielectric uncertainty dominates at $\sim 20\,\%$.}
  \label{fig:err-budget}
\end{figure}


\section{Discussion}\label{sec:disc}

The geometric absorption law has four regime boundaries that frame
its applicability.

For sources in the reactive near field ($d < \lambda/(2\pi)$, that
is $1.7\,$mm at $28\,$GHz), evanescent waves and antenna-body
impedance coupling require full-wave simulation. Outside this
regime, the law applies pointwise with spatially varying
inputs.

The local surface map $\Sab(\rr)$ loses physical meaning below
$6\,$GHz, where the skin depth exceeds the surface layer thickness
and multi-tissue Fabry-P\'erot interference grows large.
Total absorbed power stays valid via the exact flux-averaged
transmission $\Tbar$.

Body-scale Mie and resonance effects below $5\,$GHz reduce the
geometric law to a qualitative tool in that band. The Mie validation
in \cref{sec:val} quantifies the error for body-scale curvatures.

The dielectric properties of biological tissue have been measured to
within $\sim 20\,\%$ at mmWave \cite{AlekseevZiskin2007}. This
uncertainty propagates linearly into $T_0$ and dominates the error
budget at every other operating regime.


\section{Conclusion}\label{sec:conc}

The absorbed power density on biological tissue at mmWave reduces to
$\Sab(\rr) = \Sinc T_0 \Vis(\rr,\khat) \pospart{\nhat \cdot (-\khat)}$,
a closed-form expression in the surface normal, the incidence
direction, and the binary self-shadowing visibility. The collapse
rests on a pseudo-Brewster compensation between TE and TM Fresnel
transmissions that places biological tissue in the high-index regime
that Azzam \cite{Azzam2015} described for optical substrates. The
unpolarised transmission stays within $5.6\,\%$ of its
normal-incidence value $T_0$ over $[0^\circ, 75^\circ]$ at $28\,$GHz
and within $5\,\%$ root-mean-square across $0.3$ to $100\,$GHz on all
biological tissues with $|\ntilde| > 2.5$. Curvature, diffraction,
and inter-body reflections enter as bounded corrections quantified in
\cref{sec:corrections}.

The multi-source form is the forward pass of a single-hidden-layer
rectified network whose weights are the path directions and powers
from any ray tracer and whose hidden-unit gates are the per-path
visibilities, with an exact activation rather than a fit. The
dosimetry chain is differentiable end-to-end with the same arithmetic
primitives as a hardware-accelerated rendering pipeline. The
visibility factor is ambient occlusion, evaluated at frame rate on
consumer GPUs. A
$10^{12}$-cell volumetric problem reduces to $O(MN)$ surface
multiplications.

The law matches Mie theory on lossy spheres to within $5\,\%$
above $60\,$GHz and $10\,\%$ at body-part sizes in the mmWave band.
It matches full Fresnel integration on the Thelonious phantom to
$0.35\,\%$ on total power. It matches Sim4Life FDTD on the same
phantom at $7\,$GHz on the IEC/IEEE\,$63195$ peak
$4\,\mathrm{cm}^2$ metric to within $3\,\%$ direction-averaged. The
worst-case error is below the $20\,\%$ tissue dielectric
uncertainty.


\bibliographystyle{IEEEtran}
\bibliography{refs}

\end{document}
